\documentclass[12pt, a4paper]{article}
\usepackage[UTF8]{ctex}
\usepackage{amsmath}
\usepackage{amssymb}
\usepackage{geometry}
\usepackage{enumitem}
\usepackage{fancyhdr}
\usepackage{titlesec}

% 页面设置
\geometry{left=2.5cm, right=2.5cm, top=2.5cm, bottom=2.5cm}

% 页眉页脚
\pagestyle{fancy}
\fancyhf{}
\rhead{专升本数据结构习题集}
\lhead{分类整理}
\cfoot{\thepage}

% 标题格式
\titleformat{\section}{\large\bfseries}{\thesection}{1em}{}
\titleformat{\subsection}{\normalsize\bfseries}{\thesubsection}{1em}{}

\title{\textbf{第一章 概述}}
\author{}
\date{}

\begin{document}

\section*{选择题}

\begin{enumerate}[label=\arabic*.]
	\item 以下哪种数据结构中任何两个结点之间都没有逻辑关系（ ）
	      \begin{enumerate}
		      \item[A.] 树形结构
		      \item[B.] 集合
		      \item[C.] 图形结构
		      \item[D.] 线性结构
	      \end{enumerate}

	\item 要求同一逻辑结构的所有数据元素具有相同的性质，这意味着（ ）
	      \begin{enumerate}
		      \item[A.] 数据元素具有同一的特点
		      \item[B.] 不仅数据元素包含的数据项的个数相同，而且对应数据项的类型要一致
		      \item[C.] 每个数据元素都一样
		      \item[D.] 数据元素所包含的数据项的个数要相等
	      \end{enumerate}

	\item 以下说法正确的是（ ）
	      \begin{enumerate}
		      \item[A.] 数据元素是数据的最小单位
		      \item[B.] 数据项是数据的基本单位
		      \item[C.] 数据结构是带有结构的各数据项的集合
		      \item[D.] 数据结构是带有结构的数据元素的集合
	      \end{enumerate}

	\item 下面的叙述不是算法的特性的是（ ）
	      \begin{enumerate}
		      \item[A.] 有穷性
		      \item[B.] 输入性
		      \item[C.] 可行性
		      \item[D.] 可读性
	      \end{enumerate}

	\item 下面的叙述中（ ）不是设计一个好的算法应达到的目标。
	      \begin{enumerate}
		      \item[A.] 健壮性
		      \item[B.] 可读性
		      \item[C.] 高存储量
		      \item[D.] 正确性
	      \end{enumerate}

	\item 下面运算中（ ）不是数据结构所具备的基本运算。
	      \begin{enumerate}
		      \item[A.] 插入
		      \item[B.] 排序
		      \item[C.] 退出
		      \item[D.] 删除
	      \end{enumerate}

	\item 数据结构是一门研究非数值计算的程序设计问题中计算机的（ ）以及它们之间的（ ）和运算等的学科。
	      \begin{enumerate}
		      \item[A.] 数据元素 \quad B. 计算方法 \quad C. 逻辑存储 \quad D. 算法
		      \item[A.] 结构 \quad B. 关系 \quad C. 运算 \quad D. 数据映像
	      \end{enumerate}

	\item 数据结构被形式地定义为 $(K, R)$，其中 $K$ 是（ ）的有限集，$R$ 是 $K$ 上的（ ）的有限集。
	      \begin{enumerate}
		      \item[A.] 算法 \quad B. 数据元素 \quad C. 数据操作 \quad D. 逻辑结构
		      \item[A.] 操作 \quad B. 映像 \quad C. 存储 \quad D. 关系
	      \end{enumerate}

	\item 在数据结构中，从逻辑上可以把数据结构分成（ ）
	      \begin{enumerate}
		      \item[A.] 动态结构和静态结构
		      \item[B.] 紧凑结构和非紧凑结构
		      \item[C.] 线性结构和非线性结构
		      \item[D.] 内部结构和外部结构
	      \end{enumerate}

	\item 数据结构在计算机内存中的表示是指（ ）
	      \begin{enumerate}
		      \item[A.] 数据的存储结构
		      \item[B.] 数据结构
		      \item[C.] 数据的逻辑结构
		      \item[D.] 数据元素之间的关系
	      \end{enumerate}

	\item 在数据结构中，与所使用的计算机无关的是数据的（ ）结构。
	      \begin{enumerate}
		      \item[A.] 逻辑
		      \item[B.] 存储
		      \item[C.] 逻辑和存储
		      \item[D.] 物理
	      \end{enumerate}

	\item 算法分析的目的是（ ），算法分析的两个主要方面是（ ）
	      \begin{enumerate}
		      \item[(1)] A. 找出数据结构的合理性 \quad B. 研究算法中的输入和输出的关系 \quad C. 分析算法的效率以求改进 \quad D. 分析算法的易懂性和文档性
		      \item[(2)] A. 空间复杂度和时间复杂度 \quad B. 正确性和简明性 \quad C. 可读性和文档性 \quad D. 数据复杂性和程序复杂性
	      \end{enumerate}

	\item 计算机算法指的是（ ），它必须具备输入、输出和（ ）等 5 个特性。
	      \begin{enumerate}
		      \item[(1)] A. 计算方法 \quad B. 排序方法 \quad C. 解决问题的有限运算序列 \quad D. 调度方法
		      \item[(2)] A. 可行性、可移植性和可扩充性 \quad B. 可行性、确定性和有穷性 \quad C. 确定性、有穷性和稳定性 \quad D. 易读性、稳定性和安全性
	      \end{enumerate}

	\item 在以下的叙述中，正确的是（ ）
	      \begin{enumerate}
		      \item[A.] 线性表的顺序存储结构优于链式存储结构
		      \item[B.] 二维数组是其数据元素为线性表的线性表
		      \item[C.] 栈的操作方式是先进先出
		      \item[D.] 队列的操作方式是先进后出
	      \end{enumerate}

	\item 在决定选取何种存储结构时，一般不考虑（ ）
	      \begin{enumerate}
		      \item[A.] 各结点的值如何
		      \item[B.] 结点个数的多少
		      \item[C.] 对数据有哪些运算
		      \item[D.] 所用编程语言实现这种结构是否方便
	      \end{enumerate}

	\item 在存储数据时，通常不仅要存储各数据元素的值，而且还要存储（ ）
	      \begin{enumerate}
		      \item[A.] 数据的处理方法
		      \item[B.] 数据元素的类型
		      \item[C.] 数据元素之间的关系
		      \item[D.] 数据的存储方法
	      \end{enumerate}

	\item 数据的逻辑结构是由（ ）局部组成的。
	      \begin{enumerate}
		      \item[A.] 2
		      \item[B.] 3
		      \item[C.] 4
		      \item[D.] 5
	      \end{enumerate}

	\item 算法是对某一类问题求解步骤的有限序列，并具有（ ）个特性。
	      \begin{enumerate}
		      \item[A.] 3
		      \item[B.] 4
		      \item[C.] 5
		      \item[D.] 6
	      \end{enumerate}

	\item 下面说法错误的是（ ）
	      \begin{enumerate}
		      \item[(1)] 算法原地工作的含义是指不需要任何额外的辅助空间。
		      \item[(2)] 在相同的规模 $n$ 下，复杂度 $O(n)$ 的算法在时间上总是优于复杂度 $O(2^n)$ 的算法。
		      \item[(3)] 所谓时间复杂度是指最坏情况下，估计算法执行时间的一个上界。
		      \item[(4)] 同一个算法，实现语句的级别越高，执行效率越低。
		      \item[A.] (1)
		      \item[B.] (1) (2)
		      \item[C.] (1) (4)
		      \item[D.] (3)
	      \end{enumerate}

	\item （ ）是数据的逻辑结构。
	      \begin{enumerate}
		      \item[A.] 链表
		      \item[B.] 线性表
		      \item[C.] 十字链表
		      \item[D.] 顺序表
	      \end{enumerate}

	\item 数据的基本单位是（ ）
	      \begin{enumerate}
		      \item[A.] 数据元素
		      \item[B.] 记录
		      \item[C.] 数据项
		      \item[D.] 数据对象
	      \end{enumerate}

	\item 基于数据的逻辑关系，数据的逻辑结构划分为（ ）基本结构。
	      \begin{enumerate}
		      \item[A.] 4 类
		      \item[B.] 3 类
		      \item[C.] 2 类
		      \item[D.] 5 类
	      \end{enumerate}

	\item 以下数据结构中，（ ）是线性结构。
	      \begin{enumerate}
		      \item[A.] 栈
		      \item[B.] 特殊矩阵
		      \item[C.] 二维数组
		      \item[D.] 二叉树
	      \end{enumerate}

	\item 算法的空间复杂度是对算法（ ）的度量。
	      \begin{enumerate}
		      \item[A.] 空间效率
		      \item[B.] 时间效率
		      \item[C.] 健壮性
		      \item[D.] 可读性
	      \end{enumerate}

	\item （ ）不是算法具有的 5 个特性之一。
	      \begin{enumerate}
		      \item[A.] 可行性
		      \item[B.] 正确性
		      \item[C.] 有穷性
		      \item[D.] 确定性
	      \end{enumerate}

	\item 下面关于线性表的叙述中，错误的是哪一个？（ ）
	      \begin{enumerate}
		      \item[A.] 线性表采用顺序存储，必须占用一片连续的存储单元
		      \item[B.] 线性表采用顺序存储，便于进行插入和删除操作
		      \item[C.] 线性表采用链式存储，不必占用一片连续的存储单元
		      \item[D.] 线性表采用链式存储，便于进行插入和删除操作
	      \end{enumerate}

	\item 线性表是具有 n 个（ ）的有限序列。
	      \begin{enumerate}
		      \item[A.] 字符
		      \item[B.] 数据元素
		      \item[C.] 数据项
		      \item[D.] 表元素
	      \end{enumerate}

	\item 线性表采用链式存储时，其地址（ ）
	      \begin{enumerate}
		      \item[A.] 必须是连续的
		      \item[B.] 一定是不连续的
		      \item[C.] 部分地址必须是连续的
		      \item[D.] 连续与否均可以
	      \end{enumerate}
\end{enumerate}

\section*{填空题}

\begin{enumerate}
	\item 一个数据结构在计算机中的（ ）称为存储结构。
	\item 数据的逻辑结构包括（ ）、（ ）和（ ）三种结构，树形结构和图形结构合称为（ ）。
	\item 在线性结构中，第一个结点（ ）前驱结点，其余每个结点有且只有（ ）个前驱结点；最后一个结点（ ）后继结点，其余每个结点有且只有（ ）个后继结点。
	\item 在树形结构中，根结点没有（ ）结点，其余每个结点有且只有（ ）个前驱结点，叶子结点没有（ ）结点，其余每个结点的后继结点可以（ ）。
	\item 在图形结构中，每个结点的前驱结点数和后继结点数可以（ ）。
	\item 线性结构中元素之间存在（ ）关系，树形结构中元素之间存在（ ）关系，图形结构中元素之间存在（ ）关系。
	\item 算法的 5 个重要特性是（ ）、（ ）、（ ）、（ ）和（ ）。
	\item 算法可以用不同的语言描述，如果用 C 语言或 PASCAL 语言等高级语言来描述，则算法实现上就是程序了。这种说法是（ ）的。
	\item 算法效率分析可分为（ ）分析和（ ）分析。
	\item 数据结构的存储方式有（ ）、（ ）、（ ）和（ ）4 种。
	\item 数据结构包括（ ）、（ ）、（ ）三个组成部分。
	\item 数据对象是性质相同的（ ）的集合。
	\item 数据逻辑结构包括 \underline{线性结构}、\underline{树形结构} 和 \underline{图状结构} 三种类型，树形结构和图状结构合称 \underline{非线性结构}。
	\item 数据的逻辑结构分为 \underline{集合}、\underline{线性结构}、\underline{树形结构} 和 \underline{图状结构} 4 种。
	\item 衡量一个算法的优劣主要考虑正确性、可读性、健壮性和 \underline{时间复杂度} 与 \underline{空间复杂度}。
	\item 评估一个算法的优劣，通常从 \underline{时间复杂度} 和 \underline{空间复杂度} 两个方面考察。
	\item 算法的 5 个重要特性是 \underline{有穷性}、\underline{确定性}、\underline{可行性}、\underline{输入} 和 \underline{输出}。
	\item 数据结构的基本存储方法是 \underline{顺序}、\underline{链式}、\underline{索引} 和 \underline{散列} 存储。
\end{enumerate}

\section*{判断题}

\begin{enumerate}
	\item 顺序存储方式只能用于存储线性结构。（ ）
	\item 数据元素是数据的最小单位。（ ）
	\item 数据结构、数据元素、数据项在计算机中的映像（或表示）分别称为存储结构、结点、数据域。（ ）
	\item 数据的物理结构是指数据在计算机内实际的存储形式。（ ）
	\item 数据的逻辑结构是对数据元素之间关系的描述，它与数据元素之间的相对位置无关。（ ）
	\item 在决定选取何种存储结构时，一般不考虑各结点的值如何。（ ）
	\item 抽象数据类型 (ADT) 包括定义和实现两方面，其中定义是独立于实现的，定义仅给出一个 ADT 的逻辑特性，不必考虑如何在计算机中实现。（ ）
	\item 抽象数据类型与计算机内部表示和实现无关。（ ）
	\item 顺序存储方式插入和删除时效率太低，因此它不如链式存储方式好。（ ）
	\item 线性表采用链式存储结构时，结点和结点内部的存储空间可以是不连续的。（ ）
	\item 对任何数据结构链式存储结构一定优于顺序存储结构。（ ）
	\item 集合与线性表的区别在于是否按关键字排序。（ ）
	\item 线性表中每个元素都有一个直接前驱和一个直接后继。（ ）
	\item 线性表就是顺序存储的表。（ ）
	\item 取线性表的第 i 个元素的时间同 i 的大小有关。（ ）
	\item 循环链表不是线性表。（ ）
\end{enumerate}

\end{document}
